% =====================================================================
%  One Stack, Two Ledgers -- a plain-language guide to the research paper
%
%  Written for a reader who has never seen the paper and does not work in
%  finance, procurement or machine learning. Compile on Overleaf with
%  pdfLaTeX; no figures or external files are needed.
% =====================================================================
\documentclass[11pt,a4paper]{article}

\usepackage[utf8]{inputenc}
\usepackage[T1]{fontenc}
\usepackage[margin=2.5cm]{geometry}
\usepackage{enumitem}
\usepackage{booktabs}
\usepackage{xcolor}
\usepackage[hidelinks]{hyperref}
\usepackage{parskip}

\setlist{itemsep=3pt, topsep=4pt}

% a light grey box for the "in one line" summaries
\newcommand{\inoneline}[1]{%
  \medskip\noindent\colorbox{gray!10}{\parbox{\dimexpr\linewidth-2\fboxsep}{%
  \textbf{In one line:} #1}}\medskip}

\title{\textbf{One Stack, Two Ledgers}\\[4pt]
\large A plain-language guide to the research paper}
\author{Rachit Garg, Kartik Joshi, Prem Kukreja, Gagandeep Singh, Samuel Alex\\
\small SP Jain School of Global Management, Dubai}
\date{}

\begin{document}
\maketitle

\section*{What this paper is about}

This paper is about catching two kinds of financial crime with one computer
system. The first crime is \textbf{money laundering}, where criminals move
dirty money through many bank accounts to hide where it came from. The second
is \textbf{bid rigging}, where companies that are supposed to compete for a
government contract secretly agree in advance who will win.

These two crimes are normally handled by completely different people, using
completely different tools. The paper shows that they are actually the same
puzzle underneath, builds one system that works on both, and then tests that
system far more strictly than is usual. Some of the results are surprising, and
the paper reports them as they are.

\inoneline{Two crimes, one system, tested honestly.}

% ---------------------------------------------------------------------
\section{The problem}
% ---------------------------------------------------------------------

\subsection{Nothing looks wrong on its own}

Imagine you are checking bank transfers one at a time. Someone sends 40,000
dirhams to another account. Is that a crime? You cannot tell. On its own it
looks like any ordinary payment.

Now imagine twelve accounts passing that same money around in a circle over
three days. That is money laundering. But you would never see it by looking at
one transfer at a time. Each single transfer still looks normal. The crime only
appears when you step back and look at the \emph{pattern} between all of them.

Bid rigging works the same way. One company submits a bid for a government
contract. The paperwork is in order. Is that a crime? No. But if five companies
take turns winning, always price their losing bids just above the winner, and
turn out to share the same directors and the same office address, that is a
cartel. Again, no single bid gives it away. The pattern does.

\subsection{Why today's systems miss it}

Most tools used by banks and auditors check \textbf{one record at a time}. A
bank's alert system fires when one transfer crosses a threshold, say ``cash
deposit above X.'' Criminals know the thresholds, so they keep every transfer
just below them. The result is the worst of both worlds: the system floods
staff with false alarms on innocent people (industry estimates put 90 to 95
percent of bank alerts as false), while quietly missing the real networks.

Procurement auditors have it even harder. They check tenders by hand, one at a
time, on a small sample. Cartels are usually discovered only years later,
through a complaint or a confession.

There is a third difficulty: real crime is rare. In the main dataset used in the
paper, only about 2 percent of transactions are known to be illegal, and most
records have no label at all. When the thing you are looking for is that rare,
a system can be ``99 percent accurate'' by simply saying \emph{nothing is
suspicious}, which is useless.

\subsection{The key idea: two crimes, same footprints}

Here is the observation that makes this a single project instead of two. When
you draw the money flows as a map, and draw the bidding relationships as a map,
the suspicious shapes are the same.

\begin{itemize}
\item Money going around in a \textbf{circle} looks like companies
      \textbf{taking turns} winning contracts.
\item Many small deposits \textbf{funnelling into one account} looks like many
      fake losing bids \textbf{propping up one winner}.
\item One lump sum \textbf{scattered across many accounts} looks like a cartel
      \textbf{carving up a market} between its members.
\item A chain of \textbf{shell companies} with hidden common owners looks like
      ``rival'' firms that secretly \textbf{share directors}.
\end{itemize}

If the shapes are the same, then one system that learns to recognise those
shapes should work on both crimes. That is the bet the paper makes, and then
tests.

\inoneline{The crime is in the pattern, not in any single record, and the
patterns are the same in both crimes.}

% ---------------------------------------------------------------------
\section{The solution}
% ---------------------------------------------------------------------

\subsection{What the system does, step by step}

\begin{enumerate}[label=\textbf{Step \arabic*.}, leftmargin=*, itemsep=6pt]

\item \textbf{Turn everything into one map.}
Every bank account, company, and government tender becomes a dot. Every
payment, bid, or shared owner becomes a line between two dots. Whether the raw
data came from a bank or from a procurement office, it ends up in the same
format. From this point on, the system does not care which crime it is looking
at.

\item \textbf{Learn what suspicious shapes look like.}
Several different detectors examine the map. Some have been trained on past
cases where the answer is known. Others have never seen a labelled example and
simply flag anything that looks structurally odd. Each detector gives every dot
a risk score.

\item \textbf{Combine the scores fairly.}
Before the detectors' opinions are merged, each one is adjusted so that ``0.8''
means the same thing for all of them. This sounds like a small detail. It
turned out to be one of the biggest single improvements in the whole paper.

\item \textbf{Judge groups, not individuals.}
This is the heart of it. Instead of handing an investigator a list of suspicious
\emph{accounts}, the system finds tightly connected \emph{groups} of dots and
ranks the groups. A laundering ring or a cartel is a group. Flagging one member
of it is not much help; flagging the group is.

\item \textbf{Stop at a fixed budget.}
Investigators have limited hours. The system produces a short, ranked list, for
example the top 50 groups, and stops. It is built to be a triage tool for a
human, not a machine that decides on its own.

\item \textbf{Explain every flag.}
For each group on the list, the system shows which dots and lines caused the
alarm, names the shape it found (a ring, a funnel, a rotation), and points to
the specific official red-flag guideline that shape matches. An alert without
an explanation is treated as not delivered.

\item \textbf{Answer questions, but never accuse.}
An assistant lets an investigator ask things in plain language, such as ``why
was this group flagged?'' It can only read the published results, it can only
quote numbers that actually exist in the data, and it is physically prevented
from saying anyone is guilty. Every answer carries a reminder that the system
screens cases and does not judge them.
\end{enumerate}

\subsection{What the team found}

The paper is unusually frank about its results, including the ones that did not
go the team's way.

\begin{itemize}
\item \textbf{The simpler method won.} On the main bank dataset, an
      older, well-understood technique (a decision-tree method fed with a few
      network measurements) scored 0.81 on a 0-to-1 scale. The deep learning
      model, the fancier approach the project was built around, scored 0.47.
      The paper does not hide this. It confirms what a major benchmark study
      had already suggested, and it shows that the gap gets \emph{wider}, not
      narrower, when you run the experiment several times instead of once.

\item \textbf{Learning transfers between markets, but unevenly.} A model
      trained on some countries' procurement data and tested on a country it
      had never seen did better than chance in most cases. But on the largest
      country in the harder dataset, it did worse than guessing. Averaging over
      all countries would have hidden that failure; the paper reports it.

\item \textbf{Learning transfers between the two crimes in only one
      direction.} A model trained on procurement data gave useful signal when
      pointed at bank data (about 2.3 times better than chance). The reverse did
      not work. And in neither direction did transfer help until there were
      already around a thousand labelled examples available, which is exactly
      when you would need help least.

\item \textbf{The strangest result.} When the team deliberately corrupted 20
      percent of the training labels, the model's score on the final test went
      \emph{up}, from 0.48 to 0.60, while its score on the checking set used to
      pick the best model went \emph{down}. This is not a recommendation to
      corrupt your data. It is a warning that on this dataset, the standard way
      of choosing a model is looking at the wrong thing.

\item \textbf{Some cartel behaviours are nearly invisible to the map.} When
      fake cartels were planted in real procurement data, the system found
      dense clusters of colluding firms 92 percent of the time. But it almost
      never caught ``taking turns'' or fake losing bids, because those do not
      change the shape of the map, only who wins. Catching them needs a
      different kind of signal.

\item \textbf{The working system exists.} It produces 254 ranked group alerts
      for the bank data and 223 for the procurement data, each with an
      explanation attached. The assistant passed 24 of its 26 test questions
      with every number correctly grounded, and produced zero accusatory
      statements.
\end{itemize}

\subsection{What the system does not do}

It does not decide that anyone is guilty. It ranks cases for a human to look
at. It has not yet been tested with real investigators (that study is designed
but not run). And the paper openly lists a further experiment, combining both
crimes into a single map to catch people who do both, that the team has
specified in full but not carried out.

\inoneline{Map everything, judge groups not individuals, give a human a short
explained list, and report the results even when they are unflattering.}

% ---------------------------------------------------------------------
\section{The technology behind it}
% ---------------------------------------------------------------------

This section names the actual tools. Each one is explained in plain words
first.

\subsection{The data}

Five public datasets were used, no private or classified data at all.

\begin{center}
\small
\begin{tabular}{p{3.6cm}p{2.3cm}p{8.2cm}}
\toprule
\textbf{Dataset} & \textbf{Domain} & \textbf{What it is} \\
\midrule
Elliptic++ & Banking & About 200,000 real Bitcoin transactions, roughly 2\% marked illegal. The main test bed. Also used in a second form that groups transactions by the wallet behind them. \\
IBM AMLworld & Banking & Five million synthetic (computer-generated) bank transactions with every laundering scheme labelled. Used to test the system at large scale. \\
Mendeley EU cartel & Procurement & About 15,600 real contracts across seven European countries, tied to 73 proven cartel cases. The main procurement test bed. \\
Garc\'ia Rodr\'iguez & Procurement & About 64,000 real bids from six markets (Japan, Italy, Brazil, USA, Switzerland). Used to test whether learning carries from one market to another. \\
OCDS Georgia & Procurement & Real, unlabelled Georgian tender data with about 163,000 entries. Fake cartels were planted into it to see whether the system could find them. \\
\bottomrule
\end{tabular}
\end{center}

\subsection{The models}

\textbf{Graph neural networks (GNNs).} This is the deep learning at the core of
the project. An ordinary neural network looks at one item at a time. A graph
neural network looks at an item \emph{and its neighbours on the map}, and their
neighbours, and so on. That is what lets it see a ring or a funnel rather than
just a single dot. Three variants were used:
\begin{itemize}
\item \textbf{GraphSAGE}, which summarises what a dot's neighbours look like;
\item \textbf{GATv2} (graph attention), which learns to pay more attention to
      the neighbours that matter most; this was the strongest of the three on
      the banking data;
\item \textbf{R-GCN}, which handles maps where the lines mean different things
      (a payment, a bid, a shared director), used on the procurement data.
\end{itemize}
The maps were built with lines in both directions, so the model can tell
sending from receiving. Removing that single feature made the model much worse,
the biggest effect of any component the team tested.

\textbf{Focal loss.} Because real crime is rare, the model was trained with a
loss function that pays extra attention to the rare positive cases instead of
being swamped by the ordinary ones.

\textbf{Unsupervised detectors.} Two methods, \textbf{DOMINANT} and a
\textbf{graph autoencoder}, learn what a ``normal'' neighbourhood looks like and
flag anything they cannot reconstruct well. They need no labelled examples,
which is why they are the only tools that work on the unlabelled Georgian data.

\textbf{XGBoost.} A gradient-boosted decision-tree method, a mainstream
machine-learning tool that is not deep learning. It was included as the
baseline to beat. Fed with a handful of network measurements (how many
connections a dot has, how tightly knit its neighbourhood is), it beat the
GNNs on the banking data.

\subsection{From scores to a list of groups}

\textbf{Isotonic calibration} adjusts each detector's scores so that they are
comparable before combining them. \textbf{Leiden community detection} is the
algorithm that finds tightly connected groups in the map. The groups are then
ranked, near-duplicates removed, and the list cut at a fixed budget.

\subsection{Explaining a flag}

\textbf{PGExplainer} is a method that, given a GNN's decision, works out which
lines on the map mattered most to it. It was chosen over the better-known
GNNExplainer because in a head-to-head test it gave sensible answers on 49 of
50 cases against 12 of 50, in half the time. A separate \textbf{rule-based
motif matcher} then checks whether those lines form one of nine known shapes
(cycle, fan-in, fan-out, pass-through, common control, rotation, cover bid,
partition, clique) and links the shape to the matching FATF or OECD red-flag
guideline.

\subsection{The assistant}

The question-answering assistant is a \textbf{large language model} wrapped in
strict guardrails: it can only run read-only database queries, a checker rejects
any number it quotes that is not present in the underlying rows, and a
deterministic filter blocks guilt-attributing language before anything is
shown. It is released only after passing a fixed set of 26 test questions.

\subsection{How everything was tested}

The testing method is itself one of the paper's main contributions.

\begin{itemize}
\item \textbf{Strict time-based splits.} The model never sees data from the
      future of the period it is being tested on.
\item \textbf{Frozen statistics.} All normalisation is computed on training
      data only and frozen. Fixing this alone dropped the team's own headline
      number from 0.69 to 0.53, and they report the lower, correct figure.
\item \textbf{Five repeated runs.} Every deep learning result is the average of
      five runs with different random starting points, with the spread
      reported. A single lucky run is never quoted.
\item \textbf{Precision at a budget.} Instead of overall accuracy, which is
      meaningless when crime is rare, the key measure is: of the top 50 (or
      100) groups the system flagged, how many were real?
\item \textbf{Statistical significance} via 2,000 rounds of paired
      bootstrap resampling for every headline comparison.
\item \textbf{Automated leakage tests.} A dedicated test suite fails the build
      if any change lets future information leak into training.
\end{itemize}

\subsection{The software}

The pipeline is written in \textbf{Python}, with the graph models built in
\textbf{PyTorch} and \textbf{PyTorch Geometric}. Results are served through a
read-only \textbf{FastAPI} web service, and investigators use a browser-based
console. Every published number is tied to one committed configuration file,
and a repository test fails if that mapping ever drifts. The whole codebase is
covered by 391 automated tests. All experiments ran on ordinary CPUs; no GPU
was required.

% ---------------------------------------------------------------------
\section*{The whole paper in one paragraph}
% ---------------------------------------------------------------------

Money laundering and bid rigging both hide in the connections between people
rather than in any single record, and the connections form the same shapes in
both crimes. The paper builds one system that turns bank and procurement data
into a single kind of map, uses graph neural networks and other detectors to
score it, groups the results into rings and cartels, hands an investigator a
short ranked list with a plain explanation for every entry, and never
pronounces guilt. Tested strictly, the deep learning model lost to a simpler
tree-based method, learning transferred between markets unevenly and between
crimes in only one direction, and corrupting training labels oddly improved
test results. The team reports all of it, because a screening tool you cannot
trust the evaluation of is not a screening tool.

\end{document}
